\section{跨品种周期相关性与传导：以猪周期为基准}
\label{sec:corr}

本章回答本报告第一部分的核心问题：\hl{以猪周期为基准，农业各子行业之间到底有多强的
周期性关联？}为了回答这个问题，我们先固定方法口径，再分四个层次递进：
\emph{同期相关}（相关矩阵）$\to$ \emph{周期同步}（周期分量）$\to$
\emph{领先滞后}（互相关）$\to$ \emph{传导检验}（HAC 回归），
最后把结论落到\emph{股票市场}（申万子行业指数）与\emph{其他养殖品种}（牛羊、禽、水产）上。

\subsection{方法：为什么不能简单看相关系数}
\label{sec:corr-method}

农产品价格序列有很强的趋势性与结构性断点（例如 2018 年非洲猪瘟、2024 年玉米进口政策转向），
如果直接用价格\emph{水平}做相关，得到的往往是“共同的趋势”而不是“共同的周期”。
因此本章的所有计算遵循四条纪律：

\begin{enumerate}
  \item \textbf{同频率、同窗口}。所有相关系数在\emph{同一时间窗口}、\emph{同一频率}上计算，
  绝不在不同样本期上比较。主窗口取 \hl{2021 年 2 月---2026 年 8 月}
  （生猪期货 2021 年 1 月上市后的完整月份，共 67 个月度点），
  长窗口取 2015 年 2 月---2026 年 8 月（用于以生猪\emph{现货}价格指数为基准的检验）。
  \item \textbf{收益率口径为主}。基准口径是月度对数收益率
  $r_t=\ln(P_t/P_{t-1})$，并剔除换月造成的异常跳空（$|r|>35\%$ 视为拼接跳空）。
  \item \textbf{周期分量口径为辅}。用价格对 12 个月移动均线的偏离率刻画“周期位置”，
  该口径能捕捉同步性，但序列的持续性会使有效样本数远小于名义样本数，
  因此\hl{只作为定性参照，不用于显著性判断}。
  \item \textbf{显著性阈值}。主窗口 $n=67$ 时 5\% 显著性的临界值为 $|\rho|>0.24$；
  季度（3 个月）收益率口径 $n\approx 20$，临界值 $|\rho|>0.46$。
  传导回归使用 \hl{Newey--West (1987) HAC 稳健标准误，滞后阶 6}，
  以应对月度数据的自相关与异方差。
\end{enumerate}

\subsection{第一层：同期相关矩阵}
\label{sec:corr-matrix}

\begin{figure}[htbp]
\centering
\begin{tikzpicture}
\begin{axis}[
  width=11.2cm, height=10.6cm,
  xtick={0,1,...,21}, ytick={0,1,...,21},
  xticklabels={生猪,鸡蛋,玉米,淀粉,粳米,豆一,豆二,菜籽,花生,豆粕,菜粕,豆油,
               菜油,棕油,白糖,棉花,棉纱,橡胶,20号胶,苹果,红枣,纸浆},
  yticklabels={生猪,鸡蛋,玉米,淀粉,粳米,豆一,豆二,菜籽,花生,豆粕,菜粕,豆油,
               菜油,棕油,白糖,棉花,棉纱,橡胶,20号胶,苹果,红枣,纸浆},
  xticklabel style={rotate=90, anchor=east, font=\tiny},
  yticklabel style={font=\tiny},
  y dir=reverse,
  colormap={agriheat}{rgb255=(150,45,40) rgb255=(252,250,245) rgb255=(28,88,58)},
  colorbar horizontal,
  colorbar style={width=7.5cm, height=0.32cm, at={(0.5,-0.13)},
    anchor=north, xtick={-1,-0.5,0,0.5,1}, xticklabel style={font=\tiny},
    xlabel={\scriptsize 相关系数 $\rho$}, xlabel style={font=\scriptsize}},
  point meta min=-1, point meta max=1,
  axis on top, tick align=outside,
]
\addplot[matrix plot*, mesh/cols=22, point meta=explicit, shader=flat]
  table[x=i, y=j, z=rho]{data/clean/corr_m_short.dat};
\end{axis}
\end{tikzpicture}
\caption{农产品期货月度对数收益率相关矩阵（2021 年 2 月---2026 年 8 月，67 个月度观测）}
\label{fig:corr-heat}
\srcfile{新浪财经期货主力连续合约；作者计算。绿色为正相关、红色为负相关；
$|\rho|>0.24$ 为 5\% 显著性水平。}
\end{figure}

图 \ref{fig:corr-heat} 是最重要的单张图，它给出三条清晰的结论：

\begin{keybox}[相关性的三个层级]
\normalsize
\textbf{① 强相关只存在于“同一条产业链”内部。}深绿色的区块集中在两处：
\emph{油脂压榨链}（豆油---棕榈油---菜油---豆粕---豆二，$\rho$ 普遍在 $0.5\sim0.8$）
与\emph{棉花---棉纱}（同原料上下游，$\rho$ 接近 1）、\emph{玉米---玉米淀粉}
（$\rho\approx0.6$）、\emph{天然橡胶---20 号胶}。这些品种共享同一个供给来源或同一个加工链条，
因此相关性是“结构性”的、跨周期稳定的。\par\vspace{3pt}
\textbf{② 跨链条相关性普遍很弱。}养殖（生猪、鸡蛋）与粮食（玉米、粳米）之间、
养殖与软商品（白糖、棉花、橡胶）之间，$\rho$ 大多落在 $|0.1|$ 以内。\par\vspace{3pt}
\textbf{③ 与生猪相关性最高的不是玉米或豆粕，而是花生与粳米。}
这提示“同属农产品”并不构成相关性的理由，\emph{共同的驱动因子}才是。
\end{keybox}

% 由 scripts/make_tables.py 自动生成，请勿手工编辑
\begin{table}[htbp]
\centering
\small
\begin{adjustbox}{max width=\textwidth}
\begin{tabular}{@{}lrrrr@{}}
\toprule
类别对 & 品种对数 & 平均相关 & 最大相关 & 最小相关 \\
\midrule
压榨链 -- 压榨链 & 10 & 0.559 & 0.862 & 0.300 \\
软商品 -- 软商品 & 10 & 0.411 & 0.887 & 0.054 \\
油料 -- 压榨链 & 20 & 0.378 & 0.753 & 0.045 \\
粮食 -- 粮食 & 3 & 0.353 & 0.729 & 0.149 \\
粮食 -- 油料 & 12 & 0.299 & 0.573 & -0.134 \\
粮食 -- 压榨链 & 15 & 0.299 & 0.452 & 0.042 \\
油料 -- 油料 & 6 & 0.259 & 0.380 & 0.088 \\
养殖 -- 养殖 & 1 & 0.226 & 0.226 & 0.226 \\
压榨链 -- 其他 & 5 & 0.221 & 0.312 & 0.110 \\
压榨链 -- 软商品 & 25 & 0.195 & 0.565 & -0.220 \\
养殖 -- 油料 & 8 & 0.185 & 0.477 & -0.021 \\
油料 -- 其他 & 4 & 0.163 & 0.299 & 0.055 \\
养殖 -- 粮食 & 6 & 0.159 & 0.325 & -0.072 \\
油料 -- 林果 & 8 & 0.142 & 0.390 & -0.038 \\
软商品 -- 林果 & 10 & 0.134 & 0.311 & -0.027 \\
养殖 -- 压榨链 & 10 & 0.132 & 0.291 & -0.112 \\
压榨链 -- 林果 & 10 & 0.117 & 0.272 & -0.024 \\
油料 -- 软商品 & 20 & 0.114 & 0.372 & -0.139 \\
粮食 -- 软商品 & 15 & 0.081 & 0.386 & -0.176 \\
粮食 -- 林果 & 6 & 0.069 & 0.285 & -0.127 \\
林果 -- 林果 & 1 & 0.053 & 0.053 & 0.053 \\
养殖 -- 其他 & 2 & 0.052 & 0.096 & 0.009 \\
粮食 -- 其他 & 3 & 0.041 & 0.115 & -0.084 \\
养殖 -- 软商品 & 10 & 0.034 & 0.259 & -0.143 \\
养殖 -- 林果 & 4 & 0.032 & 0.127 & -0.018 \\
软商品 -- 其他 & 5 & 0.005 & 0.226 & -0.180 \\
林果 -- 其他 & 2 & -0.073 & -0.010 & -0.136 \\
\bottomrule
\end{tabular}
\end{adjustbox}
\caption{同一时间窗口（2021-02～2026-08）内大类之间的平均相关系数}
\label{tab:corr-group}
\end{table}


\begin{figure}[htbp]
\centering
\begin{tikzpicture}
\begin{axis}[agri axis, catbar,
  width=12.6cm, height=7.6cm,
  xlabel={与生猪期货月度收益的相关系数 $\rho$}, xmin=-0.3, xmax=0.55,
  yticklabels from table={data/clean/corr_with_hog.dat}{label},
  ytick=data,
  axis y line*=left,
]
\addplot[xbar, bar width=4.5pt, fill=AgriGreen!70, draw=AgriGreen]
  table[x=rho, y expr=\coordindex]{data/clean/corr_with_hog.dat};
\addplot[SoftRed, dashed, line width=0.8pt, domain=-0.3:0.55] {0};
\addplot[InkGray, dotted, line width=0.8pt, domain=-0.3:0.55] {0.24};
\addplot[InkGray, dotted, line width=0.8pt, domain=-0.3:0.55] {-0.24};
\end{axis}
\end{tikzpicture}
\caption{各期货品种与生猪期货的月度收益率相关系数（排序）}
\label{fig:corr-with-hog}
\srcfile{虚线为 0，点线为 5\% 显著性阈值 $\pm0.24$（$n=67$）。}
\end{figure}

% 由 scripts/make_tables.py 自动生成，请勿手工编辑
\begin{table}[htbp]
\centering
\small
\resizebox{\textwidth}{!}{%
\begin{tabular}{@{}llr@{}}
\toprule
品种 & 大类 & 与生猪期货的相关系数 \\
\midrule
花生 & 油料 & +0.477 \\
粳米 & 粮食 & +0.286 \\
豆一 & 油料 & +0.256 \\
鸡蛋 & 养殖 & +0.226 \\
豆油 & 压榨链 & +0.156 \\
棉纱 & 软商品 & -0.143 \\
豆粕 & 压榨链 & +0.128 \\
红枣 & 林果 & +0.127 \\
棉花 & 软商品 & -0.126 \\
菜油 & 压榨链 & +0.119 \\
菜粕 & 压榨链 & -0.112 \\
纸浆 & 其他 & +0.096 \\
玉米 & 粮食 & +0.091 \\
玉米淀粉 & 粮食 & +0.083 \\
豆二 & 油料 & +0.058 \\
天然橡胶 & 软商品 & +0.052 \\
白糖 & 软商品 & +0.026 \\
油菜籽 & 油料 & -0.021 \\
苹果 & 林果 & -0.018 \\
20号胶 & 软商品 & +0.014 \\
棕榈油 & 压榨链 & +0.003 \\
\bottomrule
\end{tabular}
}
\caption{各品种与生猪期货月度收益率的相关系数（2021-02～2026-08）}
\label{tab:corr-hog}
\srcfile{作者计算；样本 67 个月，5\% 显著性阈值 $|\rho|>0.24$。}
\end{table}


图 \ref{fig:corr-with-hog} 与表 \ref{tab:corr-hog} 给出了“以猪周期为基准”的直接答案：
\hl{没有任何一个农产品期货品种与生猪期货的月度收益率相关性达到 5\% 显著水平}。
最高的花生（$+0.48$）与粳米（$+0.29$）分别是小品种与政策定价品种，
其相关性更可能来自共同的宏观流动性因素而非产业传导；豆一（$+0.26$）
与鸡蛋（$+0.23$）处在临界值附近；而玉米（$+0.09$）、豆粕（$+0.13$）
——\emph{理论上与猪周期关联最强的两个饲料原料品种}——都不显著。

\subsection{第二层：周期同步性（周期分量口径）}
\label{sec:corr-cycle}

把价格换成“对 12 个月均线的偏离率”后，相关性结构发生了变化（图 \ref{fig:corr-cyc}）：
\hl{棉花（$-0.59$）与棉纱（$-0.66$）与猪周期呈现明显的反向同步}，
花生（$+0.52$）、纸浆（$+0.40$）同向，白糖（$-0.30$）反向，豆粕（$+0.28$）弱同向。

\begin{figure}[htbp]
\centering
\begin{tikzpicture}
\begin{axis}[agri axis, catbar,
  width=12.6cm, height=7.0cm,
  xlabel={与生猪价格周期分量的相关系数 $\rho$}, xmin=-0.75, xmax=0.6,
  yticklabels from table={data/clean/corr_cyc_hog.dat}{label},
  ytick=data,
]
\addplot[xbar, bar width=4.5pt, fill=OilBlue!70, draw=OilBlue]
  table[x=rho, y expr=\coordindex]{data/clean/corr_cyc_hog.dat};
\addplot[SoftRed, dashed, line width=0.8pt, domain=-0.75:0.6] {0};
\end{axis}
\end{tikzpicture}
\caption{各品种与生猪（现货价格指数）周期分量的相关系数}
\label{fig:corr-cyc}
\srcfile{周期分量 $=$ 月度价格 / 12 个月移动均线 $-1$；主窗口同前。
该口径序列持续性强，\emph{不用于显著性判断}。}
\end{figure}

棉花与生猪的\emph{反向}关系有清晰的经济解释：2021---2022 年是棉花大牛市
（新疆减产与全球补库）而猪价崩塌，2023---2024 年棉花回落而猪价反弹，
两者的周期错位来自完全不同的驱动因素（棉花看全球纺织需求与新疆天气，
生猪看国内产能）。
这恰恰强化了第 \ref{sec:corr-matrix} 节的结论：\hl{跨品种相关性的方向本身不稳定}，
不能作为交易或风险管理的稳定依据。

\subsection{第三层：领先滞后结构}
\label{sec:corr-leadlag}

\begin{figure}[htbp]
\centering
\begin{tikzpicture}
\begin{groupplot}[group style={group size=2 by 2, horizontal sep=1.1cm, vertical sep=1.25cm},
  agri axis, width=6.4cm, height=4.5cm,
  xlabel style={font=\scriptsize}, ylabel style={font=\scriptsize},
  tick label style={font=\tiny},
  xmin=-12, xmax=12, xtick={-12,-6,0,6,12},
  ymin=-0.9, ymax=0.9, ytick={-0.8,-0.4,0,0.4,0.8}, ylabel={$\rho$}]
\nextgroupplot[title={（a）豆粕}, xlabel={滞后月数 $k$}]
  \addplot[OilBlue, thick] table[x=lag, y=rho]{data/clean/leadlag2_M.dat};
  \addplot[SoftRed, dashed, line width=0.7pt, domain=-12:12] {0};
  \addplot[InkGray, dotted, line width=0.7pt, domain=-12:12] {0.24};
  \addplot[InkGray, dotted, line width=0.7pt, domain=-12:12] {-0.24};
\nextgroupplot[title={（b）鸡蛋}]
  \addplot[AgriGold, thick] table[x=lag, y=rho]{data/clean/leadlag2_JD.dat};
  \addplot[SoftRed, dashed, line width=0.7pt, domain=-12:12] {0};
  \addplot[InkGray, dotted, line width=0.7pt, domain=-12:12] {0.24};
  \addplot[InkGray, dotted, line width=0.7pt, domain=-12:12] {-0.24};
\nextgroupplot[title={（c）玉米淀粉}, xlabel={滞后月数 $k$}]
  \addplot[Grain, thick] table[x=lag, y=rho]{data/clean/leadlag2_CS.dat};
  \addplot[SoftRed, dashed, line width=0.7pt, domain=-12:12] {0};
  \addplot[InkGray, dotted, line width=0.7pt, domain=-12:12] {0.24};
  \addplot[InkGray, dotted, line width=0.7pt, domain=-12:12] {-0.24};
\nextgroupplot[title={（d）白糖}]
  \addplot[SoftRed, thick] table[x=lag, y=rho]{data/clean/leadlag2_SR.dat};
  \addplot[SoftRed, dashed, line width=0.7pt, domain=-12:12] {0};
  \addplot[InkGray, dotted, line width=0.7pt, domain=-12:12] {0.24};
  \addplot[InkGray, dotted, line width=0.7pt, domain=-12:12] {-0.24};
\end{groupplot}
\end{tikzpicture}
\caption{与生猪周期分量的领先滞后互相关（$k>0$ 表示该品种滞后猪周期 $k$ 个月）}
\label{fig:leadlag}
\srcfile{周期分量口径，$\pm12$ 个月，点线为 5\% 显著性阈值；仅展示相关系数绝对值较大的品种。}
\end{figure}

% 由 scripts/make_tables.py 自动生成，请勿手工编辑
\begin{table}[htbp]
\centering
\small
\resizebox{\textwidth}{!}{%
\begin{tabular}{@{}lllrrrl@{}}
\toprule
品种 & 代码 & 大类 & 同期相关 & 最大显著相关 & 滞后月 & 判定 \\
\midrule
棉花 & CF & 软商品 & -0.588 & -0.816 & 2.000 & 生猪领先 \\
棉纱 & CY & 软商品 & -0.656 & -0.804 & 2.000 & 生猪领先 \\
玉米淀粉 & CS & 粮食 & -0.167 & -0.715 & 3.000 & 生猪领先 \\
纸浆 & SP & 其他 & +0.396 & -0.679 & 9.000 & 生猪领先 \\
白糖 & SR & 软商品 & -0.305 & +0.641 & 9.000 & 生猪领先 \\
红枣 & CJ & 林果 & -0.266 & +0.613 & -9.000 & 生猪滞后 \\
苹果 & AP & 林果 & -0.417 & +0.603 & -7.000 & 生猪滞后 \\
豆二 & B & 油料 & +0.182 & -0.573 & 8.000 & 生猪领先 \\
菜粕 & RM & 压榨链 & -0.254 & -0.570 & 3.000 & 生猪领先 \\
花生 & PK & 油料 & +0.520 & +0.525 & -1.000 & 生猪滞后 \\
棕榈油 & P & 压榨链 & -0.273 & +0.521 & -5.000 & 生猪滞后 \\
豆一 & A & 油料 & -0.043 & -0.517 & 5.000 & 生猪领先 \\
油菜籽 & RS & 油料 & -0.166 & +0.470 & -6.000 & 生猪滞后 \\
豆油 & Y & 压榨链 & +0.063 & -0.458 & 8.000 & 生猪领先 \\
豆粕 & M & 压榨链 & +0.283 & -0.458 & 6.000 & 生猪领先 \\
玉米 & C & 粮食 & -0.030 & +0.362 & -4.000 & 生猪滞后 \\
菜油 & OI & 压榨链 & -0.040 & -0.341 & 5.000 & 生猪领先 \\
天然橡胶 & RU & 软商品 & -0.134 & +0.294 & -9.000 & 生猪滞后 \\
20号胶 & NR & 软商品 & -0.134 & +0.286 & -6.000 & 生猪滞后 \\
鸡蛋 & JD & 养殖 & +0.224 & \nadata & \nadata & 无显著领先滞后关系 \\
\bottomrule
\end{tabular}
}
\caption{以生猪周期分量为基准的领先滞后关系（月度，±12 个月）}
\label{tab:leadlag}
\srcfile{作者计算；比较对象为价格对 12 个月移动平均的偏离（周期分量），只报告通过 5\% 显著性的滞后阶；滞后月 $k>0$ 表示生猪领先 $k$ 个月。}
\end{table}


图 \ref{fig:leadlag} 的形状比数值更有信息量：四条互相关曲线都在
$k\in[-2,3]$ 的窄区间内出现主峰，随后迅速衰减到噪声水平。
这说明\hl{农产品价格之间的“领先滞后”关系基本上在季度以内}，
不存在可利用的、跨越半年以上的稳定传导链。对报告的结论而言，
这意味着：\emph{不能用猪价的拐点去推断 6 个月后玉米或豆粕的拐点}。

\subsection{第四层：传导检验（HAC 回归）}
\label{sec:corr-transmission}

把“相关性”升级为“条件传导”，检验模型为

\begin{equation}
  r^{i}_{t} = \alpha + \beta_{k}\, r^{\text{LH}}_{t-k} + \varepsilon_{t},
  \qquad k = 0,1,\dots,12 ,
  \label{eq:transmission}
\end{equation}

其中 $r^{\text{LH}}$ 为生猪期货月度收益率，$r^{i}$ 为目标品种收益率，
标准误为 Newey--West HAC（滞后 6 阶）。表 \ref{tab:transmission} 报告显著滞后阶数与最优滞后阶。

% 由 scripts/make_tables.py 自动生成，请勿手工编辑
\begin{table}[htbp]
\centering
\small
\resizebox{\textwidth}{!}{%
\begin{tabular}{@{}lllrrrrl@{}}
\toprule
品种 & 代码 & 大类 & 显著滞后阶数 & 最优滞后 $k$ & $\beta$ & $t$ 值 & 显著 $k$ 清单 \\
\midrule
白糖 & SR & 软商品 & 7 & 2.000 & -0.091 & -4.170 & 1,2,4,6,7,8,11 \\
棉花 & CF & 软商品 & 5 & 1.000 & -0.299 & -4.740 & 1,2,4,5,11 \\
玉米淀粉 & CS & 粮食 & 4 & 11.000 & +0.081 & 2.930 & 2,3,5,11 \\
棉纱 & CY & 软商品 & 4 & 1.000 & -0.202 & -4.390 & 1,2,4,5 \\
豆一 & A & 油料 & 3 & 2.000 & -0.129 & -3.380 & 0,2,5 \\
豆粕 & M & 压榨链 & 3 & 8.000 & -0.252 & -3.380 & 6,8,12 \\
花生 & PK & 油料 & 2 & 0.000 & +0.214 & 4.910 & 0,2 \\
20号胶 & NR & 软商品 & 2 & 2.000 & -0.106 & -2.200 & 2,7 \\
纸浆 & SP & 其他 & 2 & 8.000 & -0.199 & -2.400 & 8,9 \\
棕榈油 & P & 压榨链 & 2 & 8.000 & -0.215 & -2.530 & 2,8 \\
豆二 & B & 油料 & 2 & 8.000 & -0.288 & -3.500 & 6,8 \\
油菜籽 & RS & 油料 & 2 & 2.000 & -0.141 & -4.200 & 2,11 \\
鸡蛋 & JD & 养殖 & 1 & 0.000 & +0.191 & 2.390 & 0 \\
菜粕 & RM & 压榨链 & 1 & 12.000 & +0.193 & 2.280 & 12 \\
苹果 & AP & 林果 & 1 & 2.000 & -0.177 & -2.860 & 2 \\
豆油 & Y & 压榨链 & 1 & 8.000 & -0.213 & -3.940 & 8 \\
玉米 & C & 粮食 & 0 & \nadata & \nadata & \nadata & 无 \\
菜油 & OI & 压榨链 & 0 & \nadata & \nadata & \nadata & 无 \\
天然橡胶 & RU & 软商品 & 0 & \nadata & \nadata & \nadata & 无 \\
红枣 & CJ & 林果 & 0 & \nadata & \nadata & \nadata & 无 \\
\bottomrule
\end{tabular}
}
\caption{生猪价格向各品种的传导回归（月度收益，HAC 稳健标准误）}
\label{tab:transmission}
\srcfile{作者计算；$r^{i}_t=\alpha+\beta_k r^{\text{LH}}_{t-k}+\varepsilon_t$，Newey--West 滞后阶 6，样本 2021-02～2026-08。}
\end{table}


\begin{figure}[htbp]
\centering
\begin{tikzpicture}
\begin{groupplot}[group style={group size=2 by 1, horizontal sep=1.3cm},
  agri axis, width=6.6cm, height=4.6cm,
  xlabel style={font=\scriptsize}, ylabel style={font=\scriptsize},
  tick label style={font=\tiny},
  xmin=-0.5, xmax=12.5, xtick={0,3,6,9,12},
  ymin=-5, ymax=5, ytick={-4,-2,0,2,4}, ylabel={$t$ 值},
  legend style={font=\tiny, at={(0.02,0.03)}, anchor=south west}]
\nextgroupplot[title={（a）豆粕}]
  \addplot[OilBlue, thick, mark=*, mark size=1.2pt]
    table[x=k, y=t]{data/clean/trans_M.dat};
  \addplot[SoftRed, dashed, line width=0.7pt, domain=-0.5:12.5] {1.96};
  \addplot[SoftRed, dashed, line width=0.7pt, domain=-0.5:12.5] {-1.96};
\nextgroupplot[title={（b）鸡蛋}]
  \addplot[AgriGold, thick, mark=*, mark size=1.2pt]
    table[x=k, y=t]{data/clean/trans_JD.dat};
  \addplot[SoftRed, dashed, line width=0.7pt, domain=-0.5:12.5] {1.96};
  \addplot[SoftRed, dashed, line width=0.7pt, domain=-0.5:12.5] {-1.96};
\end{groupplot}
\end{tikzpicture}
\caption{生猪收益率对目标品种收益率的滞后传导：HAC $t$ 值随滞后阶的变化}
\label{fig:transmission}
\srcfile{模型 \eqref{eq:transmission}；虚线为 $\pm1.96$（5\% 显著性）。}
\end{figure}

结果具有高度的规律性：\hl{21 个目标品种中有 17 个至少存在一个显著滞后阶，
但显著系数都很小（$|\beta|<0.3$）且符号不一致}。
对饲料原料而言，豆粕在 $k=6,8,12$ 显著但系数为 \emph{负}（$k=8$ 时 $\beta=-0.25$，$t=-3.38$），
菜粕在 $k=12$ 显著为 \emph{正}（$\beta=+0.19$），二者方向相反、逻辑上无法同时成立，
说明这些显著性更可能来自样本内的偶然共动而非稳定的经济机制
（\emph{21 个品种 $\times$ 13 个滞后阶 $=273$ 次检验，纯随机下也会有约 14 次“显著”}）。

这一多重检验问题必须显式处理，否则会得出“存在传导”的错误结论。
本报告的处理是把显著性当作\emph{筛选}而非\emph{证明}：
只有同时满足“系数方向符合产业链逻辑”、“在季度收益率口径下仍然显著”、
“在分样本（2021---2023 / 2024---2026）中符号一致”三个条件的传导关系才被认定为稳健。
按此标准，稳健的候选只剩三个：\hl{鸡蛋（$+0.28$ / $+0.26$）}、
豆一（$+0.30$ / $+0.18$）与花生（$+0.56$ / $+0.17$），
其中只有鸡蛋有明确的产业逻辑——它与生猪同属“养殖产能周期”驱动的品种，
且是唯一在 HAC 同期回归中显著的品种（$\beta=0.19$，$t=2.39$）。
这正是第 \ref{sec:other-livestock} 节把养殖业作为独立周期组来处理的原因。

\subsection{第五层：阶段收益对比}
\label{sec:corr-phase}

相关性度量的是“同时变动”，而下文更关心“在猪价上涨/下跌阶段，其他品种表现如何”。
以生猪价格的 ZigZag 分段为状态变量，计算各品种在猪价上行段与下行段的
月均收益率，结果见图 \ref{fig:phase-scatter} 与表 \ref{tab:phase-futures}。

\begin{figure}[htbp]
\centering
\begin{tikzpicture}
\begin{axis}[agri axis,
  width=12.4cm, height=7.4cm,
  xlabel={猪价上行段：品种月均收益（\%）},
  ylabel={猪价下行段：品种月均收益（\%）},
  xmin=-4.2, xmax=2.6, ymin=-1.8, ymax=2.2,
  xtick={-4,-3,-2,-1,0,1,2}, ytick={-1.5,-1,-0.5,0,0.5,1,1.5,2},
  legend style={at={(0.99,0.03)}, anchor=south east, cells={anchor=west}},
  legend entries={各品种, 45$^\circ$ 线（两阶段同收益）},
]
\addplot[SoftRed, dashed, line width=0.9pt, domain=-4.2:2.6] {x};
\addplot[only marks, mark=*, mark size=2.2pt, AgriGreen,
  nodes near coords, point meta=explicit symbolic,
  every node near coord/.append style={font=\tiny, InkGray, anchor=west}]
  table[x=up, y=down, meta=label]{data/clean/phase_scatter.dat};
\end{axis}
\end{tikzpicture}
\caption{生猪价格上行段与下行段中，各期货品种的月均收益率}
\label{fig:phase-scatter}
\srcfile{横轴为猪价上行段（共 15 个月）内的品种月均收益，纵轴为下行段（共 47 个月）内的月均收益；
点位于 45$^\circ$ 线上方说明该品种在猪价下跌时表现更好。}
\end{figure}

\begin{figure}[htbp]
\centering
\begin{tikzpicture}
\begin{axis}[agri axis, catbar,
  width=12.4cm, height=8.2cm,
  xlabel={与猪价同向的月份占比（\%）}, xmin=35, xmax=72,
  yticklabels from table={data/clean/phase_same_fut.dat}{label},
  ytick=data,
]
\addplot[xbar, bar width=4.5pt, fill=AgriGold!75, draw=AgriGold]
  table[x=pct, y expr=\coordindex]{data/clean/phase_same_fut.dat};
\addplot[SoftRed, dashed, line width=0.8pt, domain=35:72] {50};
\end{axis}
\end{tikzpicture}
\caption{与猪价同向的月份占比（50\% 为无关系基准）}
\label{fig:phase-same}
\srcfile{同向 $=$ 品种月度收益率与生猪价格指数变动方向一致。}
\end{figure}

% 由 scripts/make_tables.py 自动生成，请勿手工编辑
\begin{table}[htbp]
\centering
\scriptsize
\begin{adjustbox}{max width=\textwidth}
\begin{tabular}{@{}lllrrrrrr@{}}
\toprule
品种 & 代码 & 大类 & 上涨段月均\% & 下跌段月均\% & 差值（pp） & 同向月占比\% & 上涨月数 & 下跌月数 \\
\midrule
生猪 & LH & 养殖 & +5.36 & -3.56 & +8.92 & \nodata & 15 & 47 \\
鸡蛋 & JD & 养殖 & +0.53 & -0.69 & +1.22 & 64.50 & 15 & 47 \\
玉米 & C & 粮食 & -0.10 & -0.32 & +0.22 & 56.50 & 15 & 47 \\
玉米淀粉 & CS & 粮食 & -0.34 & -0.23 & -0.11 & 53.20 & 15 & 47 \\
粳米 & RR & 粮食 & -0.21 & +0.02 & -0.23 & 58.10 & 15 & 47 \\
强麦 & WH & 粮食 & +0.89 & +0.78 & +0.11 & 52.40 & 7 & 14 \\
豆一 & A & 油料 & -0.69 & -0.19 & -0.50 & 59.70 & 15 & 47 \\
豆二 & B & 油料 & -0.70 & -0.06 & -0.64 & 53.20 & 15 & 47 \\
油菜籽 & RS & 油料 & -0.69 & +0.21 & -0.90 & 50.00 & 15 & 47 \\
花生 & PK & 油料 & +1.81 & -1.09 & +2.91 & 67.20 & 14 & 47 \\
豆粕 & M & 压榨链 & -0.50 & -0.23 & -0.27 & 53.20 & 15 & 47 \\
菜粕 & RM & 压榨链 & -1.79 & +0.13 & -1.92 & 54.80 & 15 & 47 \\
豆油 & Y & 压榨链 & -0.45 & +0.34 & -0.79 & 59.70 & 15 & 47 \\
菜油 & OI & 压榨链 & -0.67 & +0.18 & -0.85 & 64.50 & 15 & 47 \\
棕榈油 & P & 压榨链 & -1.26 & +1.12 & -2.38 & 53.20 & 15 & 47 \\
白糖 & SR & 软商品 & -0.28 & +0.13 & -0.41 & 48.40 & 15 & 47 \\
棉花 & CF & 软商品 & -2.84 & +0.93 & -3.76 & 43.50 & 15 & 47 \\
棉纱 & CY & 软商品 & -1.90 & +0.55 & -2.45 & 41.90 & 15 & 47 \\
天然橡胶 & RU & 软商品 & +0.10 & +0.27 & -0.17 & 43.50 & 15 & 47 \\
20号胶 & NR & 软商品 & +0.06 & +0.47 & -0.41 & 48.40 & 15 & 47 \\
苹果 & AP & 林果 & -1.85 & +1.66 & -3.51 & 45.20 & 15 & 47 \\
红枣 & CJ & 林果 & -3.34 & +0.77 & -4.12 & 58.10 & 15 & 47 \\
纸浆 & SP & 其他 & +1.06 & -0.72 & +1.79 & 50.00 & 15 & 47 \\
\bottomrule
\end{tabular}
\end{adjustbox}
\caption{生猪期货上涨段与下跌段中各品种的月均收益（2021—2026）}
\label{tab:phase-futures}
\end{table}


散点图给出了一个反直觉但极其重要的事实：\hl{大多数品种位于 45$^\circ$ 线上方}，
即\emph{它们在猪价下跌阶段的表现好于猪价上涨阶段}。
这再次印证了第 \ref{sec:grain-cycle} 节的“量与价分离”：猪价上行往往对应
养殖产能扩张（利多饲料原料的\emph{数量}），但当猪价上行由供给冲击驱动时
（如非洲猪瘟），饲料需求反而是下降的；而猪价下跌阶段常伴随成本端的原料上涨
（2021---2022 年的玉米、豆粕牛市）。同向占比最高的花生（67.2\%）与鸡蛋（64.5\%）
也印证了显著性检验的结果。

\subsection{股票市场：申万农业子行业指数}
\label{sec:corr-sw}

对第一部分而言，股票市场提供了一个额外维度的验证。用申万一级（801010 农林牧渔）
与全部二级、三级农业行业指数的月度收益率与生猪期货收益率做相关，结果见图 \ref{fig:sw-corr}。

\begin{figure}[htbp]
\centering
\begin{tikzpicture}
\begin{axis}[agri axis, catbar,
  width=12.0cm, height=6.6cm,
  xlabel={与生猪期货月度收益的相关系数 $\rho$}, xmin=-0.25, xmax=0.06,
  yticklabels from table={data/clean/sw_corr_hog.dat}{label},
  ytick=data,
]
\addplot[xbar, bar width=4.5pt, fill=HogPink!75, draw=HogPink]
  table[x=rho, y expr=\coordindex]{data/clean/sw_corr_hog.dat};
\addplot[SoftRed, dashed, line width=0.8pt, domain=-0.25:0.06] {0};
\addplot[InkGray, dotted, line width=0.8pt, domain=-0.25:0.06] {-0.24};
\end{axis}
\end{tikzpicture}
\caption{申万农业子行业指数与生猪期货的月度收益率相关系数}
\label{fig:sw-corr}
\srcfile{申万行业指数（2021 版分类）月度收益率，主窗口 67 个月度点；未达 5\% 显著性。}
\end{figure}

% 由 scripts/make_tables.py 自动生成，请勿手工编辑
\begin{table}[htbp]
\centering
\scriptsize
\begin{adjustbox}{max width=\textwidth}
\begin{tabular}{@{}lllrrr@{}}
\toprule
行业代码 & 行业名称 & 层级 & 上涨段月均\% & 下跌段月均\% & 差值（pp） \\
\midrule
801\,010 & 农林牧渔 & 一级 & -1.00 & -0.41 & -0.59 \\
801\,012 & 农产品加工 & 二级 & -3.08 & +0.26 & -3.34 \\
801\,014 & 饲料 & 二级 & -0.85 & -0.95 & +0.10 \\
801\,015 & 渔业 & 二级 & -3.76 & +1.27 & -5.03 \\
801\,016 & 种植业 & 二级 & -2.22 & +0.52 & -2.74 \\
801\,017 & 养殖业 & 二级 & +0.50 & -0.96 & +1.47 \\
801\,018 & 动物保健 & 二级 & -4.62 & +0.60 & -5.22 \\
850\,111 & 种子 & 三级 & -2.88 & +0.27 & -3.15 \\
850\,113 & 其他种植业 & 三级 & -1.62 & +0.95 & -2.57 \\
850\,122 & 水产养殖 & 三级 & -4.23 & +1.32 & -5.55 \\
850\,142 & 畜禽饲料 & 三级 & -2.90 & -0.79 & -2.11 \\
850\,151 & 果蔬加工 & 三级 & -1.35 & +1.84 & -3.19 \\
850\,152 & 粮油加工 & 三级 & -4.57 & -0.89 & -3.69 \\
850\,154 & 其他农产品加工 & 三级 & -2.33 & +0.73 & -3.07 \\
850\,172 & 生猪养殖 & 三级 & -0.28 & -0.62 & +0.35 \\
850\,173 & 肉鸡养殖 & 三级 & -1.59 & -0.41 & -1.18 \\
850\,181 & 动物保健Ⅲ & 三级 & -4.62 & +0.60 & -5.22 \\
\bottomrule
\end{tabular}
\end{adjustbox}
\caption{申万农业子行业指数在生猪上涨段与下跌段的月均收益（2021—2026）}
\label{tab:phase-sw}
\end{table}


结论非常明确：\hl{申万农业各子行业指数与猪价的月度相关性全部在 $[-0.20,+0.03]$ 之间，
无一显著}，包括理论上关联最紧密的生猪养殖（$-0.15$）与畜禽饲料（$-0.13$）。
原因是股票价格交易的是\emph{预期}\emph{与估值}，而猪价是\emph{当期实现}的价格：
在猪价上行期，市场往往已经在交易产能扩张带来的未来供给过剩（股价见顶）；
在猪价下行期，市场又开始交易产能出清带来的反转预期。
表 \ref{tab:phase-sw} 的阶段收益对比印证了这一点，并给出一个清晰的分组：

\begin{itemize}
  \item \textbf{顺周期组（仅养殖）}：\emph{养殖业指数在猪价上行段月均 $+0.50\%$、
  下行段 $-0.96\%$（差值 $+1.47$ 个百分点）}，生猪养殖指数差值 $+0.35$ 个百分点；
  这是全部申万农业行业中唯一呈现“猪价涨则强、猪价跌则弱”的行业。
  \item \textbf{反周期组（幅度最大）}：水产养殖（$-5.55$ 个百分点）、
  动物保健（$-5.22$）、渔业（$-5.03$）、粮油加工（$-3.69$）。
  这些行业在猪价下跌阶段显著跑赢——水产养殖与渔业的需求由\emph{鱼类}
  而非生猪决定，动物保健则受益于猪价低迷期的防疫投入与出栏量增长。
  \item \textbf{中性组}：饲料（$+0.10$）、农林牧渔整体（$-0.59$）、
  其他农产品加工（$-3.07$）介于两者之间。
\end{itemize}

这一分组对第二部分有一个直接含义：\hl{上市公司画像不能按“农业”整体来做，
必须按产业链位置分别处理}——同为养殖业，生猪养殖与水产养殖的收益驱动完全相反。

\subsection{本章结论}
\label{sec:corr-conclusion}

\begin{keybox}[六个可检验的结论]
\normalsize
\textbf{① 相关性有层级，没有“农业整体周期”。}强相关只存在于同一产业链内部
（压榨链、棉纺链、玉米---淀粉、橡胶双品种）；跨链条相关性几乎全部不显著。
\par\vspace{3pt}
\textbf{② 与猪周期相关性最高的品种是花生与粳米，而非玉米、豆粕。}
理论上的“饲料传导链”在价格数据上不可见。
\par\vspace{3pt}
\textbf{③ 周期分量口径下才出现的强相关（棉花 $-0.59$、棉纱 $-0.66$）方向不稳定，}
只能作为定性参考。
\par\vspace{3pt}
\textbf{④ 领先滞后关系集中在 $\pm3$ 个月以内}，不存在可用的半年以上传导链。
\par\vspace{3pt}
\textbf{⑤ 传导回归的显著性大量来自多重检验；筛选后唯一稳健的候选是鸡蛋}
（养殖产能周期同源）。
\par\vspace{3pt}
\textbf{⑥ 股票市场对猪周期几乎没有线性定价}（全部 $|\rho|<0.20$），
只有养殖业指数显示出顺周期特征。
\end{keybox}

这些结论共同指向一个判断：\hl{中国农业的周期性是“分品种的、局部的”，
而不是“整体的、同一节拍”的}。种植业、养殖业、软商品各自的周期由不同的
外生变量驱动——养殖业看产能与疫病，粮食看政策与进口，软商品看全球供需与天气，
林业果业看产地气候与库存。因此，任何“农业整体周期”的叙述都必须先说明
它指的是哪一条链、哪一个环节。
